\section*{Chapter \ref{ch:g7:stats} --- Organizing Data}

\begin{solution}{exo:g7:stats:1}
Counting each mark:
\begin{center}
\renewcommand{\arraystretch}{1.2}
\begin{tabular}{l|ccccc|c}
mark & $9$ & $11$ & $12$ & $14$ & $15$ & total \\
\hline
count & $4$ & $3$ & $7$ & $3$ & $3$ & $20$ \\
\end{tabular}
\end{center}
\end{solution}

\begin{solution}{exo:g7:stats:2}
Frequencies: $\frac{4}{20} = 20\,\%$, $\frac{3}{20} = 15\,\%$,
$\frac{7}{20} = 35\,\%$, $15\,\%$, $15\,\%$. \omterm{def:g1:addition:def}{Sum}:
$20 + 15 + 35 + 15 + 15 = 100\,\%$. \checkmark
\end{solution}

\begin{solution}{exo:g7:stats:3}
First survey: $\frac{30}{50} = 0.6 = 60\,\%$. Second:
$\frac{45}{90} = 0.5 = 50\,\%$. The first survey shows the higher
proportion, although it has fewer families with one car in absolute
count.
\end{solution}

\begin{solution}{exo:g7:stats:4}
$\intco{5}{15}$: $8, 12, 9, 14, 11$ --- count $5$.
$\intco{15}{25}$: $23, 22, 18, 16$ --- count $4$.
$\intco{25}{35}$: $31, 27, 29, 25, 33$ --- count $5$.
$\intco{35}{45}$: $35$ --- count $1$. Total $15$. \checkmark
\end{solution}

\begin{solution}{exo:g7:stats:5}
Four bars of heights $8$, $12$, $6$, $4$ over the values $0$, $1$,
$2$, $3$ (graduation every $2$ works well).
\end{solution}

\begin{solution}{exo:g7:stats:6}
Total $30$ families. Frequencies: $\frac{8}{30}$, $\frac{12}{30}$,
$\frac{6}{30}$, $\frac{4}{30}$. Angles ($\times 360^\circ$): $96^\circ$,
$144^\circ$, $72^\circ$, $48^\circ$; \omterm{def:g1:addition:def}{sum}
$96 + 144 + 72 + 48 = 360^\circ$. \checkmark
\end{solution}

\begin{solution}{exo:g7:stats:7}
Percentages: summer $\frac{162}{360} = 45\,\%$; spring
$\frac{90}{360} = 25\,\%$; autumn and winter
$\frac{54}{360} = 15\,\%$ each. Out of $60$ people, summer was chosen
by $0.45 \times 60 = 27$ of them.
\end{solution}

\begin{solution}{exo:g7:stats:8}
Class A: $\frac{12}{30} = 40\,\%$. Class B: $\frac{14}{40} = 35\,\%$.
Class A has the higher proportion of walkers, even though class B has
more walkers in count.
\end{solution}

\begin{solution}{exo:g7:stats:9}
With the axis starting at $9\,700$, the three bars have visible
heights $100$, $300$ and $400$ units: the last looks four times the
first, suggesting sales quadrupled. With an axis from $0$, the bars
are nearly equal ($9\,800$ to $10\,100$ is a rise of about $3\,\%$):
the honest picture shows almost flat sales.
\end{solution}

\begin{solution}{exo:g7:stats:10}
Frequencies sum to $1$: $f = 1 - 0.35 - 0.4 = 0.25$. Counts (total
$80$): $0.35 \times 80 = 28$; $0.4 \times 80 = 32$;
$0.25 \times 80 = 20$. Check: $28 + 32 + 20 = 80$.
\end{solution}

\begin{solution}{exo:g7:stats:11}
Girls: $55\,\%$ of $400 = 220$; boys: $180$. Cafeteria-going girls:
$0.4 \times 220 = 88$; boys: $0.6 \times 180 = 108$. Total:
$88 + 108 = 196$ students, i.e.\ $\frac{196}{400} = 49\,\%$ of the
school --- between $40\,\%$ and $60\,\%$, closer to the girls' rate
because girls are more numerous.
\end{solution}

\begin{solution}{pb:g7:stats:1}
\textbf{1.} Counts: $120 > 90$, school A wins. Frequencies:
$\frac{120}{400} = 0.30 = 30\,\%$ against
$\frac{90}{250} = 0.36 = 36\,\%$: school B wins. The prize
should follow the \omterm{def:g7:stats:frequency}{frequency} --- school B recycles a larger share
of what it uses; school A merely drinks more juice.

\textbf{2.} Teachers: $\frac{18}{30} = 0.60 = 60\,\%$. Students:
$\frac{210}{600} = 0.35 = 35\,\%$. Together:
$\frac{18 + 210}{630} = \frac{228}{630} \approx 0.36 =
36\,\%$. The students are twenty times more numerous, so the
combined \omterm{def:g7:stats:frequency}{frequency} is pulled almost entirely to their side ---
the big group carries the big weight.

\textbf{3.} Frequencies must add up to $100\,\%$
(\cref{def:g7:stats:frequency}), but
$45 + 30 + 20 + 10 = 105$: at least one sector is wrong.

\textbf{4.} Count for \omterm{def:g7:stats:frequency}{frequency} $0.15$:
$0.15 \times 40 = 6$. The counts so far: $14 + 10 + 6 = 30$, so
the last count is $40 - 30 = 10$, with \omterm{def:g7:stats:frequency}{frequency}
$\frac{10}{40} = 0.25$. (Check: $0.35 + 0.25 + 0.15 + 0.25 =
1$.)

\textbf{5.} Angles: $0.40 \times 360 = 144^\circ$,
$0.35 \times 360 = 126^\circ$, $0.25 \times 360 = 90^\circ$;
and $144 + 126 + 90 = 360^\circ$: the pie closes.

\textbf{6.} $\frac{2.4}{10} = 0.24 = 24\,\%$ of the ballots came
back.

\textbf{7.} Telephone directories and car registries listed
mostly wealthy households, whose vote differed from the
country's; the $2.4$ million answers were a giant count drawn
from the wrong population. As in question 1: what matters is
not how many you count, but whom --- a \omterm{def:g7:stats:frequency}{frequency} computed on a
distorted sample describes the sample, not the country.

\textbf{8.} True support:
$0.7 \times 1\,000 + 0.3 \times 9\,000 = 700 + 2\,700 = 3\,400$
supporters out of $10\,000$: $34\,\%$. The poll:
$0.7 \times 500 + 0.3 \times 500 = 350 + 150 = 500$ out of
$1\,000$: $50\,\%$ --- sixteen points too high, because the
wealthy are \omterm{def:g3:division:half}{half} the sample but a tenth of the town.

\textbf{9.} Replies: $0.4 \times 200 = 80$ unhappy and
$0.1 \times 800 = 80$ happy, so $160$ replies of which $80$
unhappy: measured unhappiness $\frac{80}{160} = 50\,\%$. True
unhappiness: $\frac{200}{1\,000} = 20\,\%$. Nobody lied ---
the unhappy simply reply more, and the survey hears them
louder.

\textbf{10.} True proportions: one tenth wealthy, so $100$
wealthy and $900$ modest interviews:
$0.7 \times 100 + 0.3 \times 900 = 70 + 270 = 340$ out of
$1\,000$: the repaired poll predicts $34\,\%$ --- the truth of
question 8. Fifty thousand well-chosen interviews beat two
million badly chosen ones.

\textbf{11.} Hospital A: mild $\frac{90}{100} = 90\,\%$; severe
$\frac{280}{400} = 70\,\%$; overall
$\frac{90 + 280}{500} = \frac{370}{500} = 74\,\%$.

\textbf{12.} Hospital B: mild $\frac{340}{400} = 85\,\%$;
severe $\frac{65}{100} = 65\,\%$; overall
$\frac{340 + 65}{500} = \frac{405}{500} = 81\,\%$. Hospital A
wins on mild cases ($90 > 85$) \emph{and} on severe cases
($70 > 65$).

\textbf{13.} Overall, B shows $81\,\%$ against A's $74\,\%$:
the hospital that loses in \emph{every} category wins the
total. (This reversal has a name: Simpson's paradox.)

\textbf{14.} The mixes are opposite: A's patients are mostly
severe ($400$ of $500$), B's mostly mild ($400$ of $500$).
Severe cases cure less often wherever they are treated, so A's
overall figure is dragged down by the hard cases it accepts ---
the overall is a weighted blend, and the weights differ between
hospitals. A severe patient should choose hospital A: $70\,\%$
beats $65\,\%$ in the only row that concerns them.

\textbf{15.} For instance: ``Your comparison blends mild and
severe patients, and the two hospitals treat opposite mixes ---
split by severity, hospital A cures a higher share of
\emph{both} kinds. Overall figures may only be compared when the
groups behind them are alike; otherwise the weights, not the
quality, decide the winner.''
\end{solution}
